\documentclass[conference]{IEEEtran}
\IEEEoverridecommandlockouts
% The preceding line is only needed to identify funding in the first footnote. If that is unneeded, please comment it out.
%Template version as of 6/27/2024

\usepackage{cite}
\usepackage{url}
\usepackage{amsmath,amssymb,amsfonts}
\usepackage{algorithmic}
\usepackage{graphicx}
\usepackage{textcomp}
\usepackage{xcolor}
\usepackage{moredefs,lips}
\usepackage{lmodern}
\usepackage[hidelinks]{hyperref}
\usepackage{mathtools}

\def\BibTeX{{\rm B\kern-.05em{\sc i\kern-.025em b}\kern-.08em
  T\kern-.1667em\lower.7ex\hbox{E}\kern-.125emX}}
\begin{document}

\title{
  Conference Paper Title*\\
  {\footnotesize \textsuperscript{*}Note: Sub-titles are not captured for
    https://ieeexplore.ieee.org  and should not be used}
}

\author{\IEEEauthorblockN{Ethan Brush}
  \IEEEauthorblockA{\textit{University of Illinois Urbana-Champaign} \\
    Urbana, IL, USA \\
    egbrush2@illinois.edu}
  \and \IEEEauthorblockN{Milind Kumar}
  \IEEEauthorblockA{\textit{University of Illinois Urbana-Champaign} \\
    Urbana, IL, USA \\
    milindk2@illinois.edu}
  \and \IEEEauthorblockN{Hausen Wu}
  \IEEEauthorblockA{\textit{University of Illinois Urbana-Champaign} \\
    Urbana, IL, USA \\
    hausenw2@illinois.edu}
  \and \IEEEauthorblockN{Ethan Zhang}
  \IEEEauthorblockA{\textit{University of Illinois Urbana-Champaign} \\
    Urbana, IL, USA \\
    ethanz4@illinois.edu}
}

\maketitle

\begin{abstract}
  This document is a model and instructions for \LaTeX.
  This and the IEEEtran.cls file define the components of your paper [title, text, heads, etc.].
  *CRITICAL: Do Not Use Symbols, Special Characters, Footnotes, or Math in Paper Title or Abstract.
\end{abstract}

\begin{IEEEkeywords}
  component, formatting, style, styling, insert.
\end{IEEEkeywords}

\section{Introduction}

\section{Methodology}

This study examines bus bunching on the Champaign-Urbana Mass Transit District (MTD) network and compares methods for controlling it.
The proposed approach uses MTD's public APIs to gather real-world GPS location data, which is then used to calibrate a simulation for headway distribution based on Bellei and Gkoumas~\cite{bellei2010}.
The simulation provides a common setting for testing and evaluating different control methods;
in particular, this study examines the efficacy of schedule-, headway-, prediction-, and learning-based methods~\cite{berrebi2018holding}\cite{rodriguez2023}.

\subsection{Simulating buses}

The simulation follows the \emph{transit line operation model} developed by Bellei and Gkoumas~\cite{bellei2010}.
Buses travel in one direction through an ordered sequence of stops.
At each stop, passengers alight and board before the bus departs.
Boarding is limited by bus capacity, and passengers who are not able to board remain at the stop waiting for the next bus.

For a given bus $m$ at stop $n$, the \emph{dwell time} is defined by the equation
\[
  S_{mn} \coloneq a_0 + a_1 B_{mn} + a_2 A_{mn},
\]
where $B_{mn}$ and $A_{mn}$ count boarding and alighting passengers, respectively.
The coefficients $a_0$, $a_1$, and $a_2$ represent fixed time at the stop (overhead) and the time required per boarding or alighting passenger.

The number of passengers arriving between the preceding bus's departure and the current bus's arrival are drawn from a Poisson distribution with mean equal to the product of the passenger arrival rate and the interval.
Passenger arrivals during dwell time are directly estimated using the average arrival rate; this accounts for the additional boarding time these passengers cause, and the further arrivals during that time.
On the other hand, Passengers alighting are drawn from a Binomial distribution, where each passenger alights with a stop-specific probability.

The recorded MTD vehicle location data provides travel-time estimates and dispatch headways.
These quantities are used to calibrate the paper's triangular travel-time distribution and minimum separation between successive buses.
Due to lack of recorded data, however, passenger arrival rates, alighting probabilities, and their corresponding coefficients $a_1$ and $a_2$ are assumed.

\subsection{Control Methods}

Four control methods are compared to the baseline simulation:

Schedule-based holding methods delay a bus until its scheduled departure.
Headway-based holding methods delay a bus when the time between it and its predecessor's departure is shorter than a target headway.
Prediction-based holding methods delay a bus based on estimates for subsequent buses arrival times.
The implementations of these three methods follow Berrebi et al.~\cite{berrebi2018holding}.

% NOTE: i don't know if we plan on implementing stop-skipping?
% I think it's probably best not to for consistency.
The learning-based approach is based on the reinforcement learning framework of Rodriguez et al.~\cite{rodriguez2023}\ to holding-only control, where agents use headway and passenger information to choose when and how long to delay a bus.

\subsection{Evaluating results}

% TODO: discuss experimental methodology (e.g. what are we measuring, how many trials, etc.) and discuss what visualizations we'll create & how we'll analyze them.

\section{Results}

\section{Conclusion}

\bibliographystyle{IEEEtran}
\bibliography{references}

\end{document}
